\subsection{Computational Complexity}
\label{sec:complexity}

We analyze the computational complexity of the proposed measures. All measures rely on forward reachability analysis as a fundamental building block.

\begin{prop}[Reachability Complexity]
    \label{prop:reach_complexity}
    Computing the reachable state set $S_{\text{reach}}$ for a planning
    problem $\langle P, O, I, G \rangle$ requires $O(2^{|P|} \cdot |O| \cdot |P|)$
    time and $O(2^{|P|} \cdot |P|)$ space.
\end{prop}
\begin{proof}
    The reachable state set has at most $2^{|P|}$ elements. Each iteration
    of the fixpoint computation applies up to $|O|$ operators, where each
    operator application requires $O(|P|)$ for precondition checking and
    effect computation. Each state requires $O(|P|)$ space.
\end{proof}

\begin{prop}[Unreachability Profile Complexity]
    \label{prop:ur_complexity}
    Given $S_{\text{reach}}$, computing $I^{\text{scope}}_{\text{UR}}$ and
    $I^{\text{struct}}_{\text{UR}}$ requires $O(|S_{\text{reach}}| \cdot |P|)$
    time and $O(|P|)$ additional space.
\end{prop}
\begin{proof}
    For each proposition $p \in P$, check whether $p$ appears in any
    state $s \in S_{\text{reach}}$. This requires one pass through all states.
\end{proof}

\begin{prop}[Goal Mutex Profile Complexity]
    \label{prop:mx_complexity}
    Given $S_{\text{reach}}$, computing $I^{\text{scope}}_{\text{MX}}$ and
    $I^{\text{struct}}_{\text{MX}}$ requires $O(|S_{\text{reach}}| \cdot |G|^2)$
    time and $O(|G|^2)$ additional space.
\end{prop}
\begin{proof}
    For each goal pair $(g_1, g_2)$, verify achievability of each goal
    ($O(|S_{\text{reach}}| \cdot |G|)$) and check no state contains both
    ($O(|S_{\text{reach}}| \cdot |G|^2)$). Results stored in a $|G| \times |G|$ matrix.
\end{proof}

\begin{prop}[Goal Sequencing Conflict Profile Complexity]
    \label{prop:gs_complexity}
    Computing $I^{\text{scope}}_{\text{GS}}$ and $I^{\text{struct}}_{\text{GS}}$
    requires $O(|G| \cdot 2^{|P|} \cdot |O| \cdot |P|)$ time in the worst case.
\end{prop}
\begin{proof}
    For each goal $g_1$ and each state $s \in S_{\text{reach}}$ with $g_1 \in s$,
    compute $\text{Reach}(s)$ to check if any $g_2$ is reachable.
    Each local reachability computation takes $O(2^{|P|} \cdot |O| \cdot |P|)$,
    repeated for up to $|S_{\text{reach}}|$ states per goal.
\end{proof}

Table~\ref{tab:complexity} summarizes these results.

\begin{table}[ht]
    \centering
    \caption{Computational complexity of proposed measures}
    \label{tab:complexity}
    \begin{tabular}{l|l|l}
        \hline
        Measure                                    & Time Complexity                            & Space Complexity       \\
        \hline
        $S_{\text{reach}}$                         & $O(2^{|P|} \cdot |O| \cdot |P|)$           & $O(2^{|P|} \cdot |P|)$ \\
        $I_{\text{UR}}$ (given $S_{\text{reach}}$) & $O(|S_{\text{reach}}| \cdot |P|)$          & $O(|P|)$               \\
        $I_{\text{MX}}$ (given $S_{\text{reach}}$) & $O(|S_{\text{reach}}| \cdot |G|^2)$        & $O(|G|^2)$             \\
        $I_{\text{GS}}$                            & $O(|G| \cdot 2^{|P|} \cdot |O| \cdot |P|)$ & $O(2^{|P|} \cdot |P|)$ \\
        \hline
    \end{tabular}
\end{table}

\paragraph{Practical Considerations.}

The exponential dependence on $|P|$ reflects the inherent complexity of planning: determining plan existence is PSPACE-complete~\cite{Bylander1994}. However, real planning problems often have much smaller reachable state spaces than the theoretical maximum (e.g., the Locked Door scenario has $|P| = 3$ but $|S_{\text{reach}}| = 1$). Additionally, all three measures share the $S_{\text{reach}}$ computation, amortizing this cost when computing multiple profiles. The Goal Sequencing Conflict Profile has higher complexity than the other measures due to requiring multiple local reachability computations; for expensive problems, the Unreachability and Mutex Profiles can be computed independently as partial diagnostics.

\paragraph{ASP-Based Implementation Considerations.}

The complexity analysis above assumes explicit state space enumeration. \ac{ASP}-based implementations may exhibit different practical performance characteristics due to several factors. First, modern \ac{ASP} solvers employ conflict-driven clause learning, unit propagation, and other optimizations that can prune large portions of the search space~\cite{Kuhlmann2025}. Second, symbolic representations in \ac{ASP} avoid explicit state enumeration when possible, potentially achieving sub-exponential performance on structured problems. Third, the declarative nature of \ac{ASP} encodings allows solvers to exploit problem structure automatically.

Consequently, the theoretical worst-case bounds may significantly overestimate practical computation times. The empirical evaluation (Chapter~\ref{sec:evaluation}) will assess actual scalability on benchmark problems, providing complementary evidence to the theoretical complexity analysis. Following~\cite{Kuhlmann2025}, we anticipate that \ac{ASP}-based computation will be competitive with or superior to explicit enumeration approaches for the measure classes considered here.
